\documentclass[11pt]{article}
\usepackage[margin=1.15in]{geometry}
\usepackage{hyperref}
\hypersetup{colorlinks=true, urlcolor=blue, linkcolor=black}
\setlength{\parindent}{0pt}
\setlength{\parskip}{7pt}
\pagestyle{empty}
\begin{document}

\hfill \today

\bigskip
Dr.\ Bo Li\\
Editor-in-Chief\\
IEEE Internet of Things Journal

\bigskip
Dear Dr.\ Li,

Thank you for the decision on IoT-66559-2026, ``FedKAN-IDS: Heterogeneity-Robust Federated
Kolmogorov--Arnold Networks for NetFlow-Based IoT Intrusion Detection,'' and for the
opportunity to revise. We are grateful to both reviewers, whose comments were specific
enough to be acted on rather than merely acknowledged.

The title has changed, to ``What Federated KANs Actually Buy: Learning-Rate Robustness at a
Compute Cost for NetFlow-Based IoT Intrusion Detection.'' The previous title promised
heterogeneity robustness; the revised measurements do not support that, and a title should
say which of the two a reader is being handed.

We should say at the outset that this is not the same paper with corrections. Reviewer~2's
seventh comment asked whether the learning rate had been tuned per architecture. It had
not: the submitted study used $\eta = 10^{-2}$ for every model. We therefore swept the
learning rate over six values for all four architectures on all ten seeds, and the result
changed our conclusions rather than confirming them.
The parameter-matched MLP's best learning rate is an order of magnitude away from the one
we had used, and giving it that learning rate removes most of the accuracy gap and most of
the variance gap we had reported as architectural.

Rather than restate a weakened version of the original claim, we have rewritten the paper
around what the measurements now support. The system, the datasets, the partitioning
protocol and the implementations are unchanged. The claims built on them are not, and the
revised manuscript reports smaller numbers than the original in the places where the
original was generous to our own method. Specifically: the headline binary advantage on
NF-BoT-IoT-v2 falls from $+5.49$ to $+0.76$~pp once the baseline is tuned; the
$2.6\times$ variance-reduction figure is withdrawn, since it is $0.71\times$, $0.33\times$ and
$0.87\times$ on the other three cells and becomes $1.05\times$ on NF-BoT-IoT-v2 itself once
both architectures are tuned; and the multi-class result is now reported as a range,
$+0.64$ to $+1.80$~pp depending on the summary statistic, rather than as $+1.735$~pp at
$p=0.012$ from a single round.

We also owe the Editor two corrections that neither reviewer raised and that we found while
acting on their comments.

\textbf{First, a result that did not follow.} Corollary~2 of the submitted manuscript
concluded that FedKAN is intrinsically less sensitive to client heterogeneity. It does not
follow from Theorem~1: the heterogeneity term in that bound applies identically to both
architectures, and the spline-approximation residual is an additional error term for the
KAN rather than a mechanism protecting it. The corollary asserted its conclusion instead of
deriving it. We have removed it.

\textbf{Second, a reporting choice that overstated precision.} All headline accuracies in
the submitted manuscript were taken from the final communication round. We measured the
round-to-round oscillation over the last ten rounds and it is $2.1$ to $2.7$ percentage
points across all four architectures, which is larger than one of the effects we reported
as significant. Every accuracy in the revised manuscript is the mean over the final five
rounds, and we state the estimator explicitly.

\textbf{Single-machine re-run.} Every experiment in this revision, including those that
appeared in the submitted version, was re-run on one machine. The submitted numbers came
from an RTX~4090 and a mixture of hardware would have made the new and old numbers
incomparable. We also report inference and training cost on that hardware, which speaks to
the deployment questions both reviewers raised.

\textbf{Marked-up copy.} We enclose a clean manuscript and a marked-up copy in which new
text is highlighted. We should be straightforward about one limitation. Where a passage was
rewritten rather than edited, an automatic comparison cannot always place the deleted text
in context, so parts of what was removed do not appear in the marked-up copy at all. We
measured this rather than leave it implicit: of the previous manuscript's body sentences,
$18.9\%$ are gone and $35.4\%$ of the present one is new, but the automatic comparison marks
only $54\%$ of the new sentences and $38\%$ of the deleted ones. The blue marking is reliable
where it appears; its absence does not mean a passage is unchanged. We therefore list the
substantive deletions explicitly in the response letter rather than leaving the reader to
infer them from the markup.

\textbf{Data and code.} The simulator, all model implementations, the added baselines, every
experiment reported here, and the scripts that regenerate each figure and table are
available at \url{https://github.com/haodpsut/FedKAN-IDS-v2}.

We hope the reviewers find the changes responsive. Where we have declined a suggestion we
say so and explain why rather than passing over it in silence.

\bigskip
Sincerely,

\bigskip
Phuc Hao Do, on behalf of all authors

\end{document}
